\section{Introducción}

Este manual está pensado para que \textbf{cualquier usuario autorizado} de \textbf{Machu Picchu Exclusive Tours E.I.R.L.} aprenda a usar \textbf{Odoo Online} desde cero, \textbf{sin depender del consultor} para las tareas cotidianas: registrar clientes, mover oportunidades en el embudo, cotizar, registrar pagos y llevar reservas en Proyectos.

La interfaz de su instancia está en \textbf{español (Latinoamérica)}. En este documento se citan los \textbf{nombres de menús y botones} tal como aparecen en pantalla (por ejemplo \textbf{Nuevo}, \textbf{Guardar}, \textbf{Contactos}).

El alcance cubierto aquí es: aplicaciones \textbf{CRM}, \textbf{Contactos} (o acceso a contactos desde el CRM), \textbf{Proyectos}, \textbf{Ventas} y uso de \textbf{actividades} y mensajería (chatter) en las fichas.

Como apoyo previo, puede ver un \textbf{video introductorio sobre Odoo} en YouTube (\textit{audio e interfaz en inglés}): \href{https://www.youtube.com/watch?v=nbso3NVz3p8}{\texttt{https://www.youtube.com/watch?v=nbso3NVz3p8}}.

En la parametrización vigente, el sistema usa \textbf{únicamente USD} y las cotizaciones muestran \textbf{importes sin impuestos}.

\textbf{Cómo usar este manual.} Siga el orden de los capítulos: primero acceso y pantalla principal; luego \textbf{gestión de clientes y contactos} (base de todo lo demás); después CRM, ventas y proyectos. Practique con un caso real de prospecto o viajero.
